\chapter{Phương pháp nghiên cứu}
\section{Tổng quan quy trình nghiên cứu}
Đề tài được xây dựng trên một quy trình nghiên cứu gồm tám phase được thiết kế chặt chẽ và nhất quán theo đúng tài liệu thiết kế \texttt{PIPELINE\_ViMedAQA.md}:
\begin{enumerate}
    \item \textbf{Project Setup \& Exploratory Data Analysis (EDA):} Thiết lập môi trường huấn luyện, dọn dẹp dữ liệu trùng lặp và tiến hành khảo sát thống kê mô tả.
    \item \textbf{Baseline Zero-shot (Gemini):} Thiết lập hiệu năng tham chiếu trên mô hình ngôn ngữ lớn thương mại thông qua Groq API \texttt{Llama-3.3-70B-versatile}.
    \item \textbf{Fine-tuning ViT5:} Tinh chỉnh mô hình sinh chuỗi chuyên biệt tiếng Việt ViT5-base.
    \item \textbf{Fine-tuning BARTpho:} Tinh chỉnh mô hình Sequence-to-Sequence BARTpho-word phục vụ đối sánh chéo.
    \item \textbf{Unified Evaluation:} Đánh giá đồng bộ hiệu năng của tất cả các mô hình trên cùng tập test và cùng bộ độ đo chuẩn hóa.
    \item \textbf{Error Analysis \& Ablation Study:} Phân tích định tính các lỗi tiêu biểu của mô hình và thực hiện nghiên cứu ablation study trên các chiến lược giải mã.
    \item \textbf{Model Publishing:} Đóng gói và đẩy trọng số mô hình lên HuggingFace Hub.
    \item \textbf{Gradio Web RAG Interface:} Xây dựng ứng dụng hỏi đáp kết hợp tìm kiếm ngữ cảnh y học và triển khai demo lên HuggingFace Spaces.
\end{enumerate}

\section{Kiến trúc sinh câu trả lời dựa trên Retrieval-Augmented Generation (RAG)}
Để giải quyết triệt để hai thách thức lớn nhất của các mô hình ngôn ngữ nhỏ (Small Language Models - SLMs) trong miền y khoa là giới hạn tri thức nền (parametric memory) và hiện tượng ảo giác thông tin (factual hallucination), đề tài nghiên cứu và triển khai hệ thống hỏi đáp theo kiến trúc tinh chỉnh \textbf{RAG (Retrieval-Augmented Generation)}. Kiến trúc đề xuất là một pipeline đa tầng, được kiểm duyệt thông qua 3 cấu phần chính:

\subsection{Cấu phần truy xuất song song (Hybrid Retrieval)}
Khi nhận được một câu hỏi y tế $q$ từ người dùng, hệ thống sẽ chuẩn hóa (chuyển chữ thường, xóa khoảng trắng thừa) và tiến hành tìm kiếm trên tập tri thức y học độc lập (\texttt{medical\_corpus.json}). Nhằm tận dụng tối đa thế mạnh của các phương pháp tìm kiếm, hệ thống kích hoạt đồng thời hai cơ chế truy xuất:
\begin{itemize}
    \item \textbf{Truy xuất thưa (Sparse Retrieval - BM25Okapi):} Câu hỏi được tách từ tiếng Việt qua thư viện \texttt{PyViTokenizer} và đưa vào thuật toán BM25Okapi để tính điểm tương quan dựa trên tần suất từ khóa. Phương pháp này xuất sắc trong việc bắt chính xác các từ khóa "cứng" như tên bệnh lý hoặc tên biệt dược. Kết quả trả về top $50$ tài liệu tốt nhất.
    \item \textbf{Truy xuất dày đặc (Dense Retrieval - Bi-Encoder + FAISS):} Hệ thống sử dụng mô hình biểu diễn ngữ nghĩa \texttt{bkai-foundation-models/vietnamese-bi-encoder} để ánh xạ câu hỏi vào không gian vector $384$ chiều. Vector này được chuẩn hóa $L_2$ và tìm kiếm độ tương đồng trong không gian \textbf{FAISS IndexFlatIP} (sử dụng Inner Product). Phương pháp này giúp nắm bắt ý nghĩa khái niệm sâu sắc ngay cả khi từ ngữ bề mặt không trùng khớp. Kết quả cũng trả về top $50$ tài liệu có khoảng cách gần nhất.
\end{itemize}

\subsection{Cấu phần lai ghép và tái xếp hạng (Fusion \& Re-ranking)}
Để tổng hợp và đánh giá lại kết quả từ hai bộ tìm kiếm độc lập, hệ thống áp dụng luồng xử lý hai bước nhằm chắt lọc ra các ngữ cảnh chính xác nhất:
\begin{enumerate}
    \item \textbf{Lai ghép thứ hạng (Reciprocal Rank Fusion - RRF):} Các tài liệu ứng viên từ BM25 và FAISS được gộp lại và tính điểm trọng số chéo. Hệ thống thiết lập trọng số ưu tiên ngữ nghĩa cao hơn ($w_{faiss} = 0.65$, $w_{bm25} = 0.35$) với hằng số làm mượt $K = 50$. Điểm RRF của mỗi tài liệu $d$ được tính theo công thức:
    $$RRF\_Score(d) = \frac{0.35}{50 + rank_{bm25}(d)} + \frac{0.65}{50 + rank_{faiss}(d)}$$
    Thuật toán tiến hành sắp xếp và chọn ra top $50$ tài liệu tiềm năng nhất từ danh sách lai ghép.
    
    \item \textbf{Tái xếp hạng và lọc nhiễu (Cross-Encoder):} Top $50$ ứng viên tiếp tục được đưa qua mô hình \texttt{mmarco-mMiniLMv2-L12-H384-v1}. Khác với Bi-Encoder, Cross-Encoder tương tác chéo trực tiếp cặp (Câu hỏi, Ngữ cảnh) để sinh ra giá trị thô (logit). Giá trị này được biến đổi qua hàm kích hoạt Sigmoid ($\sigma$) để tính toán độ tin cậy thực tế:
    $$prob\_score = \sigma(logit) = \frac{1}{1 + e^{-logit}}$$
    Hệ thống thiết lập một chốt chặn an toàn với ngưỡng $\tau = 0.4$. Các ngữ cảnh có $prob\_score < \tau$ sẽ bị đánh dấu là dữ liệu nhiễu và lập tức bị loại bỏ. Hệ thống sẽ lấy top-$K$ (mặc định $K=3$) ngữ cảnh vượt qua kiểm duyệt để chuyển sang bộ sinh.
\end{enumerate}

\subsection{Cấu phần sinh câu trả lời (Generator)}
Nếu không có ngữ cảnh nào vượt qua màng lọc kiểm duyệt, hệ thống kích hoạt cơ chế Fallback, tự động từ chối trả lời để đảm bảo an toàn y khoa. Ngược lại, tập ngữ cảnh sạch thu được sẽ được ghép nối với câu hỏi ban đầu $q$ theo cấu trúc Prompt nghiêm ngặt:
$$X = \text{``question: } q \text{ context: } c_1 \ . \ c_2 \ . \ \dots \ c_K\text{''}$$

\noindent Chuỗi $X$ được SentencePiece Tokenizer phân tách và đưa vào mô hình Sequence-to-Sequence \texttt{ViT5-base} (\texttt{ntthanh0307/vit5-vimedaq-medical-qa}). Quá trình giải mã tự hồi quy sử dụng thuật toán Beam Search với cấu hình tối ưu hiệu năng (số chùm \texttt{num\_beams=2}, hệ số phạt lặp \texttt{repetition\_penalty=1.5}, hệ số phạt độ dài \texttt{length\_penalty=2.0}). Các tham số này ép mô hình phải bám sát ngữ liệu đầu vào để sinh ra câu trả lời y tế minh bạch, chính xác và giảm thiểu rủi ro sinh từ vựng ảo giác.

\subsection{Sơ đồ kiến trúc tổng thể}
Dưới đây là sơ đồ kiến trúc tổng thể của dự án:

\begin{figure}[H]
    \centering
    \resizebox{\textwidth}{!}{% Đảm bảo co giãn đồng bộ toàn bộ sơ đồ vừa khít trang A4
\begin{tikzpicture}[
    font=\sffamily\small,
    >=Stealth,
    node distance=0.8cm and 0.8cm,
    % --- Styles cho các node con ---
    box/.style={rectangle, draw=black!60, thick, rounded corners, align=center, minimum width=3.8cm, minimum height=0.9cm, fill=white},
    db/.style={cylinder, draw=gray!80, thick, aspect=0.25, shape border rotate=90, align=center, fill=gray!10, minimum width=3.2cm, minimum height=1.2cm},
    diamondbox/.style={diamond, draw=red!60, thick, fill=white, aspect=2, inner sep=2pt, align=center},
    % --- Styles màu cho từng Khối ---
    col1/.style={box, draw=blue!60, fill=blue!5},
    col2/.style={box, draw=orange!60, fill=orange!5},
    col3/.style={box, draw=green!60!black, fill=green!5},
    col4/.style={box, draw=purple!60, fill=purple!5},
    title/.style={font=\bfseries, align=center}
]

% ==========================================
% KHỐI 1: PIPELINE XỬ LÝ DỮ LIỆU (ETL)
% ==========================================
\node[db] (c1_raw) {Dữ liệu y tế thô\\(VimedQA Raw)};
\node[col1, below=of c1_raw] (c1_clean) {Lọc trùng \&\\Làm sạch văn bản};
\node[col1, below=of c1_clean] (c1_chunk) {Chunking theo câu\\(Sentence-aware)};
\node[col1, below=of c1_chunk] (c1_meta) {Làm giàu Metadata\\([Keyword] Title)};
\node[col1, below=of c1_meta] (c1_json) {Tệp Corpus JSON\\(Sẵn sàng Index)};

\draw[->, thick] (c1_raw) -- (c1_clean);
\draw[->, thick] (c1_clean) -- (c1_chunk);
\draw[->, thick] (c1_chunk) -- (c1_meta);
\draw[->, thick] (c1_meta) -- (c1_json);

% ==========================================
% KHỐI 2: ĐÁNH CHỈ MỤC ĐA NHÁNH
% ==========================================
\node[col2, right=2.5cm of c1_clean] (c2_dense_model) {Bi-Encoder Model\\(vietnamese-sbert)};
\node[col2, below=of c2_dense_model] (c2_faiss) {FAISS Vector DB\\(Nhánh Dense)};

\node[col2, below=1.5cm of c2_faiss] (c2_sparse_model) {PyVi Tokenizer\\(Tách từ TV)};
\node[col2, below=of c2_sparse_model] (c2_bm25) {BM25 Index\\(Nhánh Sparse)};

\draw[->, thick] (c1_json.east) -- ++(1cm,0) |- (c2_dense_model.west);
\draw[->, thick] (c1_json.east) -- ++(1cm,0) |- (c2_sparse_model.west);

\draw[->, thick] (c2_dense_model) -- (c2_faiss);
\draw[->, thick] (c2_sparse_model) -- (c2_bm25);

% ==========================================
% KHO LƯU TRỮ CỤC BỘ (HUGGING FACE)
% ==========================================
\node[db, right=1.8cm of c2_faiss, yshift=-2cm, fill=teal!10, draw=teal!80] (hf_storage) {Tệp tin cục bộ\\(Hugging Face Space)};

% Biểu diễn thao tác Upload file thủ công/Git LFS
\draw[->, thick, dashed] (c2_faiss.east) -| node[midway,above,xshift=-1.5cm,font=\scriptsize\ttfamily] {faiss\_index.bin} (hf_storage.north);
\draw[->, thick, dashed] (c2_bm25.east) -| node[below, xshift=-1.5cm,font=\scriptsize\ttfamily] {bm25\_index.pkl} (hf_storage.south);

% ==========================================
% KHỐI 3: LUỒNG RAG CỐT LÕI (RUN-TIME)
% ==========================================
\node[col3, right=6.5cm of c2_dense_model, yshift=1.5cm] (c3_user) {Người dùng\\(Gửi câu hỏi)};
\node[col3, below=of c3_user] (c3_prep) {Chuẩn hóa \&\\Từ đồng nghĩa};
\node[col3, below=of c3_prep] (c3_hybrid) {Truy xuất song song\\(Hybrid Search)};
\node[col3, below=of c3_hybrid] (c3_rrf) {Hợp nhất \&\\Nội suy điểm (RRF)};
\node[col3, below=of c3_rrf] (c3_cross) {Cross-Encoder\\(Chấm điểm \& Rerank)};
\node[diamondbox, below=0.5cm of c3_cross] (c3_thresh) {Score > 0.4?};

\draw[->, thick] (c3_user) -- (c3_prep);
\draw[->, thick] (c3_prep) -- (c3_hybrid);

% Ứng dụng đọc file tĩnh trực tiếp từ ổ cứng của Hugging Face
\draw[->,thick] (hf_storage.east)--(c3_hybrid.west);

\draw[->, thick] (c3_hybrid) -- (c3_rrf);
\draw[->, thick] (c3_rrf) -- (c3_cross);
\draw[->, thick] (c3_cross) -- (c3_thresh);

% ==========================================
% KHỐI 4: TỔNG HỢP & HIỂN THỊ UI
% ==========================================
% �� Đà LOẠI BỎ PROMPT ENGINEERING, NỐI ĐẠT TRỰC TIẾP VÀO ViT5
\node[col4, right=2cm of c3_hybrid] (c4_vit5) {ViT5 Seq2Seq\\(Tổng hợp câu trả lời)};
\node[col4, below=1cm of c4_vit5] (c4_out) {Giao diện Gradio\\(Hiển thị kết quả)};
\node[col4, below=1cm of c4_out] (c4_fallback) {Fallback\\(Cảnh báo từ chối)};

% Điều hướng từ Diamond (Pass đi vào LLM, Fall đi vào Fallback)
\draw[->, thick] (c3_thresh.east) -- ++(1,0) |- node[above, xshift=0.5cm, font=\bfseries] {Đạt} (c4_vit5.west);
\draw[->, thick] (c3_thresh.south)--++(0,-0.45cm)--++(5.9cm,0)node[midway,above, font=\bfseries, xshift=-0.2cm] {Thấp} --++(0,2cm)--(c4_fallback.south);

\draw[->, thick] (c4_vit5) -- (c4_out);
\draw[->, thick, dashed] (c4_fallback) |- (c4_out.south);


%\draw[->, thick, dashed, draw=blue!50] (c4_out.east) -- ++(-1.5cm,0) |- (c3_user.west);
% Thêm nhãn cho đường phản hồi nếu cần: node[near start,above,rotate=90] {Gửi câu hỏi tiếp}

% ==========================================
% VẼ KHUNG BACKGROUND CHO CÁC KHỐI
% ==========================================
\begin{scope}[on background layer]
    % Background Khối 1
    \node[rectangle, draw=blue!80, fill=blue!5, thick, rounded corners, inner sep=15pt, fit=(c1_raw) (c1_json)] (bg1) {};
    \node[title, text=blue!80!black, anchor=south] at (bg1.north) {KHỐI 1: PIPELINE XỬ LÝ\\DỮ LIỆU (ETL)};
    
    % Background Khối 2
    \node[rectangle, draw=orange!80, fill=orange!5, thick, rounded corners, inner sep=15pt, fit=(c2_dense_model) (c2_bm25)] (bg2) {};
    \node[title, text=orange!80!black, anchor=south] at (bg2.north) {KHỐI 2: ĐÁNH CHỈ MỤC\\ĐA NHÁNH (INDEXING)};
    
    % Background Khối 3
    \node[rectangle, draw=green!80!black, fill=green!5, thick, rounded corners, inner sep=15pt, fit=(c3_user) (c3_thresh)] (bg3) {};
    \node[title, text=green!60!black, anchor=south] at (bg3.north) {KHỐI 3: LUỒNG RAG\\CỐT LÕI (RUN-TIME)};
    

    \node[rectangle, draw=purple!80, fill=purple!5, thick, rounded corners, inner sep=15pt, fit=(c4_vit5) (c4_fallback) (c4_out)] (bg4) {};
    \node[title, text=purple!80!black, anchor=south] at (bg4.north) {KHỐI 4: TỔNG HỢP\\\& HIỂN THỊ UI};
\end{scope}

\end{tikzpicture}
    }
    \caption{Sơ đồ kiến trúc tổng thể của hệ thống ViMedAQA}
    \label{fig:architecture_diagram}
\end{figure}

\section{Quy trình huấn luyện và tối ưu hóa Seq2Seq}
Quá trình tinh chỉnh (fine-tuning) hai mô hình chính ViT5-base \cite{phan2022vit5} và BARTpho-word \cite{tran2021bartpho} sử dụng hàm mất mát âm log-likelihood (Negative Log-Likelihood - NLL) trên từng token của chuỗi đáp án mục tiêu:
\begin{equation}
  \mathcal{L}(\theta) = -\sum_{t=1}^{T}\log p_\theta(a_t \mid a_{<t}, X)
\end{equation}
Trong đó $\theta$ đại diện cho các trọng số mạng nơ-ron được tối ưu hóa.

Để đảm bảo quá trình đối sánh chéo diễn ra công bằng, chúng tôi thiết kế quy trình huấn luyện đồng bộ về mặt kích thước batch hiệu dụng (Effective Batch Size), thuật toán tối ưu hóa và số lượng epoch huấn luyện. Tuy nhiên, do kiến trúc tokenizer và dung lượng bộ nhớ của hai mô hình có sự khác biệt lớn, các siêu tham số kỹ thuật được tinh chỉnh linh hoạt nhằm tránh lỗi tràn bộ nhớ (Out of Memory - OOM) trên GPU Nvidia T4 ($16$GB VRAM) của Google Colab. Bảng~\ref{tab:training_configs} tổng hợp chi tiết cấu hình huấn luyện thực tế của các mô hình.

\begin{table}[H]
\centering
\caption{Cấu hình siêu tham số huấn luyện thực tế trong đề tài}
\label{tab:training_configs}
\resizebox{\textwidth}{!}{%
\begin{tabular}{lcc}
\toprule
\textbf{Siêu tham số} & \textbf{Mô hình chính (ViT5-base)} & \textbf{Mô hình đối sánh (BARTpho-word)} \\
\midrule
Kiến trúc cơ sở & T5 (Text-to-Text) \cite{raffel2020t5} & BART (Seq2Seq Autoencoder) \cite{lewis2020bart} \\
Dung lượng mô hình & 220M parameters & 228M parameters \\
Kích thước Batch huấn luyện & per-device batch size = 4 & per-device batch size = 2 \\
Gradient Accumulation Steps & 4 steps & 8 steps \\
\textbf{Batch size hiệu dụng} & \textbf{16} & \textbf{16} \\
Tối ưu hóa (Optimizer) & AdamW ($\beta_1=0.9, \beta_2=0.999$) & AdamW ($\beta_1=0.9, \beta_2=0.999$) \\
Tốc độ học (Learning Rate) & $3 \times 10^{-4}$ & $3 \times 10^{-4}$ \\
Lịch trình giảm tốc độ học & Linear Warmup ($10\%$ steps) & Linear Warmup ($10\%$ steps) \\
Số Epoch huấn luyện & 5 epochs & 5 epochs \\
Độ dài đầu vào cực đại & 512 tokens & 1024 tokens \\
Độ dài đầu ra cực đại & 128 tokens & 256 tokens \\
Kiểu dữ liệu số thực & mixed precision (\texttt{fp16}) & mixed precision (\texttt{fp16}) \\
Gradient Checkpointing & Tắt & \textbf{Bật (Bắt buộc để tránh OOM)} \\
\bottomrule
\end{tabular}%
}
\end{table}

Việc duy trì kích thước batch hiệu dụng bằng $16$ giúp hướng gradient di chuyển ổn định trong không gian tham số, hỗ trợ cả hai mô hình hội tụ tốt. Cơ chế \textbf{Gradient Checkpointing} được kích hoạt bắt buộc cho BARTpho-word để đánh đổi tài nguyên tính toán lấy dung lượng bộ nhớ, cho phép mô hình xử lý chuỗi đầu vào dài lên tới $1024$ tokens mà không gặp sự cố phần cứng. Quá trình huấn luyện cũng tích hợp cơ chế lưu checkpoint tự động và lưu nhật ký huấn luyện định kỳ lên Google Drive để phòng ngừa rủi ro mất kết nối đột ngột trong môi trường Colab. Checkpoint cuối cùng được lựa chọn dựa trên chỉ số tổn thất kiểm định (\texttt{eval\_loss}) thấp nhất trên tập validation.